\section*{Chapter \ref{ch:ml:text-from-bag-of-words-to-embeddings} --- Text: From Bag of Words to Embeddings}

\begin{solution}{exo:ml:text-from-bag-of-words-to-embeddings:1}
Tokens: \texttt{<co>}, did, not, cut, guidance. Bigrams: \texttt{<co> did}, did not, not cut, cut guidance. ``Not cut''
carries the negation; the word ``cut'' alone reads as bad news.
\end{solution}

\begin{solution}{exo:ml:text-from-bag-of-words-to-embeddings:2}
$\log(20\,000/3\,000) = 1.90$ and $\log(20\,000/200) = 4.61$: ``halves'' weighs about 2.4 times as much. Rare words are
more informative about which document this is, which is not the same as more informative about the return.
\end{solution}

\begin{solution}{exo:ml:text-from-bag-of-words-to-embeddings:3}
$(1 - 1)/6 = 0$ for the compound (``beat'' positive, ``cut'' negative), planted $-1.2\%$; $1/4 = 0.25$ for the award
(``wins'' positive), planted 0.
\end{solution}

\begin{solution}{exo:ml:text-from-bag-of-words-to-embeddings:4}
They appear in the same slot with the same neighbours (a company, then ``estimates'', ``forecasts'' or ``expectations''), and
co-occurrence is all the embedding sees: cosine 0.95. Only a signal that differs between them separates them: returns (a
supervised fine-tune), or contexts that differ (``beat \dots\ shares rise''), which a larger corpus supplies in part.
\end{solution}

\begin{solution}{exo:ml:text-from-bag-of-words-to-embeddings:5}
Not a bug here: in this corpus ``to be'' appears only in takeover headlines, so it carries their return. It becomes one
when the phrasing that made it informative changes (a new wire's style, a new template), or when the word stands for
something outside the text, such as one outlet's coverage of one kind of company; a screened word should be read, and
the list checked for stability across folds.
\end{solution}

\begin{solution}{exo:ml:text-from-bag-of-words-to-embeddings:6}
Today's symbol file maps old tickers to nothing and reused tickers to today's holder; the errors concentrate in the past,
in companies that changed or disappeared, which is exactly where survivorship and takeover effects live. A hand check of
recent links, or of links that resolved, measures the wrong population. Here today's aliases give 97.5\% precision among
linked mentions but only 87.5\% correct overall, and 91.5\% precision on tickers. Check a sample stratified by year and by
mention kind, including the unlinked ones, against point-in-time aliases.
\end{solution}

\begin{solution}{exo:ml:text-from-bag-of-words-to-embeddings:7}
Unmasked, the tf-idf model's IC falls from 0.273 to 0.266: the company words get small weights fitted to noise, since no
company effect was planted. The supervised screen keeps no company word and scores 0.238 either way, because its
three-standard-error bar rejects them. On real data company words would carry each company's returns over the training
period (a good year, a takeover), a leak of past performance dressed as text, which masking prevents.
\end{solution}

\begin{solution}{exo:ml:text-from-bag-of-words-to-embeddings:8}
For a linear model with intercept $b$ and word weights $w$, the compound ``$A$ but $C$'' scores $b + w_A + w_{\text{but}} + w_C
= f(A) + f(C) + (w_{\text{but}} - b)$, where $f(A)$ and $f(C)$ are the scores of the single-clause headlines with the same
words. Matching $f(\text{beat}) = 1$, $f(\text{cut}) = -1.5$ and $f(\text{beat but cut}) = -1.2$ requires $w_{\text{but}} - b =
-0.7$; matching $f(\text{missed}) = -1$, $f(\text{raised}) = 1.5$ and $f(\text{missed but raised}) = 1.2$ requires $+0.7$. The
two conditions contradict each other. Bigrams such as ``but cut'' and ``but raised'', or an interaction between the clause
after ``but'' and a guidance indicator, can fit both: the second clause dominates.
\end{solution}

\begin{solution}{pb:ml:text-from-bag-of-words-to-embeddings:1}
\textbf{Part I.}
\begin{enumerate}
\item So that no model learns a company's identity (its past returns in the sample) as a text signal, and so that one model
serves every company.
\item Gives it weight $\log(N/N) = 0$: it cannot distinguish documents.
\item From 123 words to 473 features, nearly four times as many; the compounds they could capture are 9\% of the headlines,
and the added variance cost more than the fit gained.
\item The planted effect explains all that text can: an IC of about 0.3, because the reaction's noise (3\%) is large against
the effects (up to 2\%). Any score above that on a \omterm{def:ml:why-financial-machine-learning-is-different:sets}{test set} is luck or leakage.
\end{enumerate}
\textbf{Part II.}
\begin{enumerate}[resume]
\item Event rules 0.292, tf-idf words 0.273, tf-idf with bigrams 0.257, supervised words 0.238, dictionary 0.224,
embeddings 0.094.
\item The tf-idf word model passes the dictionary between 1\,000 and 3\,000 labelled headlines; the supervised list and
the bigram model only beyond 3\,000.
\item They were written from the corpus's own templates and read all 20\,000 headlines correctly; they fail on a phrasing
they were not written for (``surpasses consensus'' returns nothing). Measure them on text written after the rules.
\item The embeddings learned from co-occurrence put opposites together (``beat'' and ``missed'' at 0.95), so the average of
a headline's vectors barely distinguishes good from bad news.
\end{enumerate}
\textbf{Part III.}
\begin{enumerate}[resume]
\item 99.5\%; its only errors are the short name Apex with no sector word, where the most covered Apex is chosen.
\item 87.5\%. Old names and tickers resolve to nothing (89.7\% of mentions linked), and reused tickers resolve to the new
holder (precision 91.5\% on tickers).
\item The wrong link: an unlinked headline is a missed trade, a wrongly linked one is a trade in the wrong security, with a
return that looks like noise or, worse, like a signal.
\item 100\% of 737 with a sector word, 58.0\% of 245 without.
\end{enumerate}
\textbf{Part IV.}
\begin{enumerate}[resume]
\item \emph{Which Apex?} Point-in-time linking is right on 99.5\% of the mentions (precision 99.5\%); linking with today's
aliases is right on 87.5\% (precision 97.5\% overall, 91.5\% on tickers). The return-supervised models reach an IC of 0.273
(tf-idf words, \omterm{def:ml:linear-and-regularised-baselines:logistic}{logistic regression}) and 0.238 (the supervised word list) against the dictionary's 0.224.
\item Linking: every model downstream inherits its errors, and its errors are silent.
\item Sample links stratified by year and mention kind, including the unlinked, and compare with point-in-time aliases;
count links to delisted or renamed companies.
\item Fit and select it in walk-forward folds on earlier text, fix the reaction window in advance, and re-screen on a
schedule with the list's stability checked.
\item Organising long documents and building features for a supervised model; not measuring tone.
\item Tokens per document rise to thousands, dictionaries become more stable, section structure and changes from the
previous filing carry information, and a mention may concern a peer, a supplier or a customer rather than the filer.
\item Few labels, a new domain, or a need for a transparent score; and lists built for the domain.
\item Getting the security right.
\end{enumerate}
\end{solution}

\begin{solution}{iq:ml:text-from-bag-of-words-to-embeddings:1}
Term counts weighted by the log inverse share of documents containing the term, rows normalised. It shrinks common words
towards zero and puts documents of different lengths on one scale, so a penalised linear model spends its budget on
distinctive words.
\iqlookfor{the formula, the reason for the idf, and length normalisation.}
\end{solution}

\begin{solution}{iq:ml:text-from-bag-of-words-to-embeddings:2}
General lists call common financial words negative (liability, tax, cost, capital) and miss financial meanings; Loughran
and McDonald (2011) found that almost three-quarters of the negative words of a widely used general dictionary are not
negative in 10-Ks, and built finance lists.
\iqlookfor{misclassification of domain words, with the reference or an example.}
\end{solution}

\begin{solution}{iq:ml:text-from-bag-of-words-to-embeddings:3}
Screen words by their association with subsequent returns on training data, weigh them, score documents; or fit a
penalised \omterm{def:ml:linear-and-regularised-baselines:logistic}{logistic regression} on tf-idf. Leaks: returns before the text's timestamp (use the time the text became
available), company names standing for their returns, labels from overlapping windows, and vocabulary chosen on the full
sample.
\iqlookfor{a supervised construction and at least two concrete leaks with remedies.}
\end{solution}

\begin{solution}{iq:ml:text-from-bag-of-words-to-embeddings:4}
Embeddings are trained to reflect shared contexts, and a word and its opposite fill the same contexts. So distances in the
space measure topic and syntax, not polarity; as features they need a supervised layer or \omterm{def:ml:representation-learning:ssl}{fine-tuning} on the target.
\iqlookfor{the distributional reason and the consequence for sentiment.}
\end{solution}

\begin{solution}{iq:ml:text-from-bag-of-words-to-embeddings:5}
Recognise mentions (gazetteer and a tagger), resolve candidates from point-in-time aliases and tickers, disambiguate by
context, output a permanent identifier with a confidence and abstain below a threshold; log everything. Failure modes:
ticker reuse, renames, shared short names, subsidiaries and brands, and a symbol file refreshed without history.
\iqlookfor{point-in-time resolution, disambiguation, abstention, and named failure modes.}
\end{solution}

\begin{solution}{iq:ml:text-from-bag-of-words-to-embeddings:6}
First the plumbing: timestamps (text time against trade time), linking changes, a vendor's format or coverage change. Then
leakage in the backtest: vocabulary or labels chosen on the full sample, the reaction window. Then drift: new phrasings,
crowding of the signal (Book~7, chapter~28), and finally the expected shrinkage of an overfitted model.
\iqlookfor{an ordering from data problems to leakage to genuine decay.}
\end{solution}
